% =============================================================================
% manuscript.tex — Ecological Informatics (Elsevier), elsarticle class
%
% GENERATED FILE. Do not edit by hand: it is produced from the Markdown sources
% by paper/tex/build_tex.mjs. Edit the Markdown and re-run the script, or the
% next build will silently discard your changes.
%
% Build:  pdflatex manuscript && bibtex manuscript && pdflatex manuscript x2
% Needs:  elsarticle.cls and elsarticle-harv.bst (TeX Live / MiKTeX / Overleaf)
%
% COMPILED. First compiled 2026-08-14 with MiKTeX 25.12 (pdfTeX), clean: no
% errors, no undefined references, no undefined citations. The first compile was
% a debugging pass and found four real defects, all fixed in this generator: a
% true minus inside a code span (the Unicode table was applied to prose but not
% to code spans), non-ASCII inside a verbatim block, figure paths resolved from
% the wrong directory, and l-columns that cannot wrap, which ran the widest
% table 763pt past the margin. See build_report.md.
% =============================================================================
\documentclass[preprint,review,12pt]{elsarticle}

\usepackage[utf8]{inputenc}
\usepackage[T1]{fontenc}
\usepackage{lmodern}
\usepackage{textcomp}
\usepackage{tabularx}
\graphicspath{{../}{../figures/}{./}{./figures/}}
\usepackage{graphicx}
\usepackage{amsmath}
\usepackage{booktabs}
\usepackage{longtable}
\usepackage{array}
\usepackage{url}
\usepackage{lineno}
\modulolinenumbers[5]

\journal{Environmental Modelling and Software}

\begin{document}

\begin{frontmatter}

%% Title. POSITIONING.md lists three candidates and directs that under Branch A
%% (Evia enlarged and re-run cleanly, which is what happened) the third is used.
%% Referee round 2, 2026-08-14: that third candidate began "Marginal diagnostics
%% miss conditional failures", which asserts more than Section 5.4 concedes. On
%% ten effective pairs the null diagnostics are "not shown to order transfer",
%% not "shown not to". The title now says what Section 4.4 found and what
%% Section 5.4 defends. Candidate 1 was not used because "do not transfer
%% between them" is contradicted by the paper's own matrix, and candidate 2 was
%% not used because "transferability cost" asserts a debit measured only by the
%% feature-removal ablation. Change here if that decision is revisited.
%% Retitled 2026-08-14 in the split, to carry the three findings the paper was
%% cut to. "Sign reversal" names the mechanism, which the old title left out;
%% "do not order" is kept verbatim because it is the only claim Section 5.4
%% defends, the diagnostics having been shown not to order transfer here rather
%% than shown incapable of ordering it anywhere.
\title{Decided before the model: a preprocessing uncertainty budget for satellite fire-susceptibility modelling across five Mediterranean regions}

%% Author block. Affiliation and corresponding address supplied by the authors
%% 2026-08-14. Still optional and not supplied: department or faculty within the
%% university, and ORCIDs.
\author[inst1]{Yunus Emre Cogurcu\corref{cor1}}
\ead{ycogurcu@cu.edu.tr}
\author[inst1]{Emrehan Metin}
\cortext[cor1]{Corresponding author.}
\affiliation[inst1]{organization={\c{C}ukurova University},
                    city={Adana},
                    country={T\"urkiye}}

\begin{abstract}
Satellite fire-susceptibility models are published with a cross-validated skill figure. Between the satellite and that figure sits a chain of preprocessing decisions, each defensible, few reported in enough detail to repeat and almost none reported with a number attached to what they cost.

This paper attaches those numbers. One pipeline, five Mediterranean wildfire regions, one label product and one classifier are held fixed while eight preprocessing decisions are varied in turn, so that their effects are comparable with each other and with the +0.056 to +0.153 ROC-AUC effect the models themselves report.

Two decisions turn out to be predictable in advance at no cost. Whether the choice of compositing chain can move a result at all is decided by scenes per distinct acquisition date. It is above one in three of five regions and exactly 1.0 in the other two, and where it is 1.0 the two chains produce identical rasters and every boundary estimate is exactly zero. The count follows the area's WRS geometry: every region spanning two rows exceeds one scene per date, and every region spanning one row does not. Whether sea contaminates a scene-fitted dryness index is decided by the area's water fraction and by where the modelled population sits in the index's binning variable; here the contamination is real, with dry edges of 28.8 to 29.9 \textdegree{}C in a marine area against 49 \textdegree{}C inland, and does not reach the modelled population.

Two provenance splits leave no trace in the outputs. The coarse thermal input is quality-screened in two regions and unscreened in three, split by export date rather than by design. The induced change is not a uniform offset: it correlates with elevation at +0.615 in one region. Propagating that region's screening through the whole downstream chain changes 92 \% of its downscaled cells, by up to 10.9 \textdegree{}C, and moves no reported association by more than +0.0003. The same three regions encode sea as exactly 0.0 \textdegree{}C, at 8.1 \% and 38.3 \% of pixels in two of them and not at all in the third.

The analysis cell is also not the cell it is called: built by aggregating a fixed count of reference pixels in a geographic coordinate system, it is 510 m north to south and 390 to 407 m east to west, 17 to 20 \% smaller in area than the cell it approximates, so spatial blocking is weaker in longitude than every block-size label suggests.

Four axes were tested and did not threaten the result, including label omission, which corrects in the direction opposite to the convenient one. The within-region results rebuild bit for bit on a different operating system, while cross-region estimates carry a $\pm$0.02 to 0.03 tolerance across library versions. We close with six items that cost nothing to report and would let a reader judge most of this budget without repeating any of it.
\end{abstract}

%% No highlights file yet.

\begin{keyword}
preprocessing uncertainty \sep reproducibility \sep sensitivity analysis \sep remote sensing pipelines \sep wildfire susceptibility \sep data provenance
\end{keyword}

\end{frontmatter}

\linenumbers

\section{Introduction}\label{sec:1}

Satellite fire-susceptibility models are published with a cross-validated skill figure and a description of the classifier. Between the satellite and that figure sits a chain of decisions: which pixels to accept, how to composite them, how to normalise an index, how to aggregate to an analysis cell, which cells to admit, and which library to run it all in. Each is defensible. Few are reported in enough detail to be repeated, and almost none are reported with a number attached to what they cost.

This paper attaches those numbers. It takes one pipeline, five Mediterranean wildfire regions, one label product and one classifier, holds all of them fixed, and varies eight preprocessing decisions in turn. The result is a budget rather than a caveat list: the effects are measured on the same data and are therefore comparable with each other, and with the effect size the models themselves report.

\subsection{Why a budget and not a caveat list}\label{sec:1.1}

A caveat list says that a choice could matter. A budget says how much, on data the reader can see. The difference matters because preprocessing decisions are not equally consequential, and treating them as a uniform fog of uncertainty is as unhelpful as ignoring them. Two of the eight axes here move a reported association by more than the interval width that the same paper's headline result carries. Two others cannot move it at all, for reasons that are visible in the inputs before any model is fitted. Knowing which is which is worth more than knowing that preprocessing is hard.

\subsection{What makes a preprocessing finding generalisable}\label{sec:1.2}

The obvious objection to an audit of one pipeline is that its findings are that pipeline's bugs. We take that seriously and answer it in three ways.

First, wherever possible an axis is converted into a \textbf{rule stated in terms of the inputs}, not of our outputs. Whether the compositing choice can move a result is decided by how many scenes share an acquisition date, which any user can count before running anything. Whether sea water contaminates a scene-fitted dryness index is decided by the AOI's water fraction and by where the modelled population sits in the index's binning variable.

Second, the failure modes we document are structural rather than incidental. A quality-screening rule added to an export script partway through a project splits a cohort \textbf{by export date rather than by design}, and nothing in the outputs records that. An analysis grid built by aggregating a fixed number of reference pixels in a geographic coordinate system is square in degrees and not on the ground. Neither is specific to this code.

Third, the axes that did \textbf{not} move the result are reported alongside those that did. A budget in which everything matters is not a budget.

\subsection{Relation to the companion paper}\label{sec:1.3}

The companion paper asks whether a fitted relationship between pre-fire thermal dryness and burning transfers between regions, and reports that it does not, that the diagnostics available before deployment do not order the failure, and that the mechanism is a reversal in the sign of the predictor-burning association. This paper shares that cohort and pipeline and asks a different question: what the observations on both sides of that comparison are made of. Where an axis here bears on a claim there, it is noted, and in every case tested the companion paper's finding survives the variation. Neither paper's conclusion depends on the other's.

\subsection{Contributions}\label{sec:1.4}

\begin{enumerate}
  \item \textbf{A measured preprocessing budget} over eight axes on five regions, with data, model and validation held fixed, so the axes are comparable with each other and with the reported effect.
  \item \textbf{Two decisions whose relevance is predictable at no cost}, from the scene inventory and from the land-cover and index distributions respectively.
  \item \textbf{Two provenance splits that leave no trace in the outputs}: a quality-screening rule applied to part of a cohort by export date, whose induced change correlates with terrain rather than being a uniform offset, and a nodata convention that encodes sea as a physical temperature of exactly 0.0 \textdegree{}C in the same regions.
  \item \textbf{A quantified statement of what the analysis cell actually is}, with the derivation reproducing every frozen cell count in the cohort exactly.
  \item \textbf{Four axes tested and found not to threaten the result}, one of which corrects in the direction opposite to the convenient one.
  \item \textbf{A reproducibility floor}: bit-identical rebuild of the within-region results on a different operating system from released code and archived data, against a $\pm$0.02 to 0.03 tolerance across library versions for the cross-region estimates.
\end{enumerate}

Section~\ref{sec:2} places the work against the reproducibility and sensitivity literature. Section~\ref{sec:3} describes the cohort and the protocol by which each axis is varied. Section~\ref{sec:4} reports the budget axis by axis. Section~\ref{sec:5} assembles it and states what should be reported as a result. Section~\ref{sec:6} concludes.



\section{Related work}\label{sec:2}

Sensitivity of environmental model results to preprocessing is not a new concern, but it is usually raised one decision at a time and inside a study whose purpose is something else. The closest established practice is in spatial validation, where random cross-validation over autocorrelated cells is known to inflate skill estimates \citep{Roberts2017,Ploton2020} and blocked designs are the standard remedy \citep{Valavi2019,Meyer2018}, itself contested as introducing a pessimistic bias \citep{Wadoux2021,Mila2022,deBruin2022}. That literature is a budget for one axis, and it is the model this paper follows for eight.

For the observations themselves, the relevant reference points are product validation rather than pipeline sensitivity: the burned-area product's omission and commission characteristics \citep{Boschetti2019,Giglio2018} bound anything built on it, and land surface temperature products carry their own documented calibration and quality behaviour \citep{Cook2014,Malakar2018}. What is rarely reported is what a given handling of those products costs in the units of the study's own result, which is the gap this paper addresses.

The transferability question that motivates the cohort is treated in the companion paper and in the studies it engages \citep{Dimarco2026,Liu2025,Meyer2021,Ludwig2023}. \citet{Xu2026} argue that wildfire modelling conclusions depend strongly on evaluation design and task formulation; the present paper is the same argument carried one layer further down, into how the observations were made.



\section{Cohort and protocol}\label{sec:3}

\subsection{The cohort and what is held fixed}\label{sec:3.1}

Five Mediterranean wildfire regions are used: Manavgat 2021 and Muğla 2021 in Türkiye, Bejís 2022 in Spain, North Evia 2021 in Greece and Montiferru 2021 in Sardinia. They were assembled for the companion study and are described there in full: place-based rectangular areas of interest, burned labels from MODIS MCD64A1 Collection 6.1, an analysis grid of $\sim$510 m cells reconstructed from a 30 m reference grid, predictors from Landsat 8 Collection 2 Level-2 and MODIS MOD11A1 over a predictor window closing the day before the first labelled burning, and a natural-vegetation primary population.

Everything downstream of the axis under test is held fixed: the same random forest with the same hyperparameters and seed, the same nested baseline and thermal feature sets, the same five-fold spatially blocked cross-validation, and the same spatial-block bootstrap with 1000 replicates. This is what makes the axes comparable: any difference reported below is the axis, not a model choice.

\subsection{How an axis is varied}\label{sec:3.2}

Each axis is varied by re-running the pipeline's own code with one input or one setting changed, and never by reimplementing it. Where a rule had to be applied that the released tooling does not expose, the pipeline's own primitives were imported and driven directly, so that the binning, percentile, alignment and metric functions are identical between arms and only the pixel selection or the parameter differs.

Three practices are used throughout and are what make the comparisons trustworthy rather than merely plausible.

\textbf{Reproduce before you vary.} Every arm that has a frozen counterpart reproduces it first, and the reproduction is reported alongside the new number. Where an arm cannot reproduce its comparator, the comparison is not made.

\textbf{Vary one thing.} Where an axis required staging inputs into a different location, the provenance pointers were repointed to the staged copies and the change recorded, rather than allowing a path mismatch to be read as a scientific difference.

\textbf{Report the axes that did nothing.} An axis found not to move the result is reported at the same length as one that did, because a budget in which every entry is large is not a budget.

\subsection{The eight axes}\label{sec:3.3}

\begin{enumerate}
  \item \textbf{Analysis-cell geometry.} Derived from the export scale and coordinate system and checked against the frozen cell counts. Not a variation but a description; reported because block-size labels depend on it.
  \item \textbf{Coarse-thermal quality screening.} The pipeline's own \texttt{QC\_\allowbreak{}Day} bit rule and three-observation minimum, applied to the three regions exported before the rule existed, compared against their unscreened means at the input.
  \item \textbf{Nodata convention.} Exact-zero counts in each region's coarse thermal layer, against the pipeline's own suspicious-zero guard and each region's water-dominant cell share.
  \item \textbf{Compositing chain.} The released counterfactual audit, which builds a date-balanced composite alongside the production scene-weighted one and compares them at support boundaries with a paired bootstrap, run for four regions. Each region's diagnostic chain is checked against the frozen canonical products before its comparison is read.
  \item \textbf{Index normalisation.} The dryness index's wet and dry edges refitted on land pixels alone, and a third set fitted once over the pooled land pixels of all five regions, with the index recomputed and aggregated to cells each time and the signed association re-measured under the registered 10-cell spatial-block bootstrap.
  \item \textbf{Label quality.} Three probes: restriction to burned cells above sub-pixel agreement thresholds; relabelling against an independent active-fire observation; and the exposure left where the pre-label safeguard did not run, measured against the unclipped burned-area product.
  \item \textbf{Derived channels and design choices.} Restriction to low gap-fill cells; removal of the two coordinate-bearing channels; the water content of the secondary population; matching the largest region to the smallest on cell count and positive count; and the blocking scale of the transfer intervals.
  \item \textbf{Software.} Library-version sensitivity from frozen probes, and an independent rebuild of the environment on a different operating system with a field-by-field comparison against the frozen metrics.
\end{enumerate}

\subsection{Reporting conventions}\label{sec:3.4}

Signed univariate associations are reported as raw ROC-AUC against the burned label and are never folded to max(AUC, 1 $-$ AUC), so a value below 0.5 is a direction rather than weakness, and a difference between regions is called supported only when the two bootstrap intervals are disjoint. Where a verdict changes on a margin smaller than the precision of the counts behind it, that is stated with the verdict.



\section{The budget, axis by axis}\label{sec:4}

Each axis below is varied with everything else held fixed: the same five regions, the same label product, the same feature sets, the same classifier, the same spatially blocked folds and the same bootstrap. Effects are therefore comparable across axes. For scale, the within-region effect these models report is +0.056 to +0.153 ROC-AUC, and its bootstrap interval at 1 km blocking is about $\pm$0.010 wide.

\subsection{The analysis cell is not the cell it is called}\label{sec:4.1}

The grid is built by aggregating 17 $\times$ 17 reference pixels, each exported at 30 m in EPSG:4326, into a cell described throughout as $\sim$500 m. Because the export is in a geographic coordinate system, the reference pixel is a step in degrees, 30 / 111,319.49 = 0.00026949\textdegree{}, so the cell is 17 times that, 0.0045814\textdegree{}, \textbf{in both axes}. On the ground it is about 510 m north to south everywhere, but only 407 m east to west at 37\textdegree{} N and 390 m at 40\textdegree{} N. Cell area is 0.199 to 0.208 km$^{2}$ against the MODIS cell's 0.250 km$^{2}$, so the analysis cell is 17 to 20 \% smaller than the cell it approximates.

The derivation is checkable rather than nominal: dividing each AOI's span by 0.0045814\textdegree{} and rounding outward reproduces every frozen cell count in the cohort exactly (Table~\ref{tab:1}).

\begin{table}[htbp]
  \centering
  \small
  \caption{The analysis cell's true ground dimensions, per region. Columns and rows are derived by dividing each AOI's span by the cell step of 0.0045814\textdegree{} and rounding outward; the frozen cell count is the one recorded in the pipeline's own outputs, and the two agree exactly in every region. The north-south edge is about 510 m everywhere; only the east-west edge varies, with latitude.}
  \label{tab:1}
  \begin{tabularx}{\linewidth}{l>{\raggedright\arraybackslash}Xrr}
    \hline
    Region & Derived columns $\times$ rows & Frozen cell count & East-west edge \\
    \hline
    Manavgat 2021 & 175 $\times$ 138 & 24,150 & 407 m \\
    Bejís 2022 & 153 $\times$ 103 & 15,759 & 391 m \\
    Muğla 2021 & 393 $\times$ 186 & 73,098 & 407 m \\
    North Evia 2021 & 175 $\times$ 131 & 22,925 & 397 m \\
    Montiferru 2021 & 66 $\times$ 49 & 3,234 & 390 m \\
    \hline
  \end{tabularx}
\end{table}

Two consequences follow and neither is cosmetic. Every block-size label in a study of this kind is the north-south dimension: a 10-cell block described as $\sim$5 km is 5.1 km by 3.9 to 4.1 km, so \textbf{spatial blocking is systematically weaker in longitude by about a fifth}, in analyses whose central robustness argument is blocking scale. And because the cell is smaller than a MODIS cell and not aligned to it, a typical cell draws on more than one MODIS cell, which dilates the labelled burned footprint relative to the source product.

\subsection{Observation provenance: a quality rule applied to part of a cohort}\label{sec:4.2}

The coarse thermal input is \texttt{MOD11A1}. Two of the five regions apply a \texttt{QC\_\allowbreak{}Day} bit rule with a three-observation minimum and write an explicit $-$9999 nodata sentinel; three apply no quality mask and declare no nodata value. The split follows \textbf{export date}, not design: the rule was added to the export script on 2026-07-23, after three regions had already been exported. Nothing in the outputs marks the difference; it is visible only by comparing metadata across regions.

The size and shape of what that costs were measured by recomputing each unscreened region's predictor-window mean with and without the rule (Table~\ref{tab:2}).

\begin{table}[htbp]
  \centering
  \small
  \caption{Cost of the quality-screening rule in the three regions that were exported without it. Each region's predictor-window mean was recomputed with the \texttt{QC\_\allowbreak{}Day} bit rule and the three-observation minimum applied, and differenced against the frozen unscreened mean. The final column is what matters: a shift correlated with elevation is not an offset a downscaling model would absorb.}
  \label{tab:2}
  \begin{tabularx}{\linewidth}{lrr>{\raggedleft\arraybackslash}X>{\raggedleft\arraybackslash}X}
    \hline
    Region & Mean shift & SD of shift & Pixels below the 3-observation minimum & Correlation of the shift with elevation \\
    \hline
    Manavgat 2021 & +1.17 \textdegree{}C & 0.60 & 0.2 \% & \textbf{+0.615} \\
    Muğla 2021 & +0.67 \textdegree{}C & 0.50 & 0.1 \% & +0.443 \\
    Bejís 2022 & +0.47 \textdegree{}C & 0.89 & 3.6 \% & $-$0.041 \\
    \hline
  \end{tabularx}
\end{table}

Coverage is barely affected. What matters is that \textbf{the change is not a uniform offset}, which a downscaling model fitted to Landsat would largely absorb. In two of the three regions it correlates with elevation, at +0.615 in the region whose elevation-burning association is the one that reverses in the companion paper. The candidate therefore has a mechanism, and it is strongest where the anomaly is.

\textbf{It does not survive propagation.} The measurement above stops at the input, so Manavgat's entire downstream chain was rebuilt from a quality-screened MODIS input and compared against an unscreened arm rebuilt the same way. Both arms were rebuilt here rather than compared against the frozen numbers, so that the contrast isolates the screening rather than mixing it with rebuild drift; the unscreened arm reproduces the published increment, at [+0.055, +0.079] against [+0.055, +0.078].

The screening propagates substantially. It changes the downscaled surface on \textbf{22,304 of 24,150 cells}, by a mean of 0.365 \textdegree{}C and a maximum of 10.9 \textdegree{}C. The modelled associations do not follow. No signed univariate association moves by more than \textbf{+0.0003}, elevation's stays at 0.374 [0.290, 0.472] in both arms, and the population and the within-region increment are unchanged.

The reason is structural rather than fortunate. Elevation is a DEM variable that screening a thermal product cannot touch, and here the screening did not change which cells are valid either. The thermal channels that carry the signal are Landsat-derived, and fusion falls back on the MODIS-derived surface across only 2.14 points of coverage, so even a 10.9 \textdegree{}C change is diluted to +0.0003 in the association. The elevation correlation at the input is real. It does not reach the model. Details are in \texttt{paper/\allowbreak{}modis\_\allowbreak{}qc\_\allowbreak{}downstream\_\allowbreak{}propagation.\allowbreak{}md}.

\textbf{Zero-fill.} The same three unscreened regions declare no nodata value, so cells with nothing to report are written as exact 0.0 \textdegree{}C. Counting them: 518 of 6,390 pixels in Manavgat (8.11 \%), \textbf{7,347 of 19,190 in Muğla (38.29 \%)} and \textbf{none at all in Bejís}. Both non-zero fractions exceed the pipeline's own 5 \% suspicious-zero guard. What the zeros are is settled by the counts, which track each AOI's water-dominant cell share almost exactly, 8.3 \% and 38.9 \% against 0.1 \% for the one inland AOI: \textbf{the zero-fill is sea encoded as a physical temperature}, not missing observation.

\subsection{Compositing: a decision whose relevance is readable in advance}\label{sec:4.3}

A current-period composite is a median over the clear acquisitions in the window, and how those acquisitions are weighted is a choice. A scene-weighted median counts every scene; a date-balanced median counts every calendar date once. The pipeline offers both, and an audit compares them at support boundaries, where a positive value means the date-balanced chain lowers the discontinuity.

For Manavgat the intervention helps at every boundary type, which is what motivated it. Extending the audit to the other four regions turns that into a rule (Table~\ref{tab:3}).

\begin{table}[htbp]
  \centering
  \scriptsize
  \caption{Whether the compositing choice can act, read from the scene inventory alone. Scenes per distinct acquisition date is counted before anything is fitted. Where it is 1.0 the two compositing chains produce identical rasters and there are no same-day boundaries to compare, so the verdict is not "no evidence" but "no effect is possible". The verdict column reports the boundary-discontinuity audit only for the three regions where the intervention can act.}
  \label{tab:3}
  \begin{tabularx}{\linewidth}{l>{\raggedright\arraybackslash}Xrrrrrr>{\raggedright\arraybackslash}X}
    \hline
    Region & WRS tiles in the window & Paths & Rows & Scenes & Dates & Scenes per date & Same-day edges & Verdict \\
    \hline
    Manavgat 2021 & 177/34, 177/35, 178/34, 178/35 & 2 & \textbf{2} & 14 & 7 & \textbf{2.0} & 172 & \textbf{supported reduction} \\
    Muğla 2021 & 179/34, 179/35, 180/34, 180/35, 181/34 & 3 & \textbf{2} & 18 & 11 & \textbf{1.6} & 473 & uncertain \\
    Bejís 2022 & 198/32, 198/33, 199/32, 199/33 & 2 & \textbf{2} & 16 & 8 & \textbf{2.0} & 79 & uncertain \\
    North Evia 2021 & 183/33, 184/33 & 2 & 1 & 8 & 8 & 1.0 & \textbf{0} & no effect \\
    Montiferru 2021 & 193/32 & 1 & 1 & 4 & 4 & 1.0 & \textbf{0} & no effect \\
    \hline
  \end{tabularx}
\end{table}

Where scenes per date is 1.0 the two chains produce \textbf{identical rasters}. Every boundary estimate is exactly zero, and no same-day boundaries exist to compare.

The dividing line is not the number of Landsat paths an AOI spans. Evia spans two paths and shows nothing, because those two paths image it on different days. Muğla spans \textbf{three}, more than any other region here, and it does show an effect. What the three responding regions share is spanning two WRS \textbf{rows}. A single overpass then delivers two scenes bearing the same date, which is exactly what a scene-weighted median double-counts. Row span, not path span, is the readable predictor.

Three regions can respond, and only Manavgat improves consistently. Bejís improves at the boundary the intervention targets, at +0.378 [+0.278, +0.483], and gets \textbf{worse} at unique-date-count edges, at $-$0.067 [$-$0.095, $-$0.041]. Muğla shows the same shape and the largest single reduction in the cohort: +0.195 [+0.172, +0.221] at the same-day multiplicity edge, against $-$0.031 [$-$0.046, $-$0.016] at unique-date-count edges. Both carry an overall verdict of uncertain.

The downstream ROC-AUC comparison is admissible only for Manavgat, where three defensible chains on an identical cohort give increments of +0.045, +0.064 and +0.084, a tolerance of about $\pm$0.02. For the other four regions the released tool declines that comparison, either because the seam verdict is uncertain or because the canonical reproduction could not be resolved.

\textbf{The rule.} Scenes per distinct acquisition date, counted from the scene inventory before anything is fitted, tells a reader whether this decision can move their result at all. In this cohort it can in \textbf{three regions of five}, and does consistently in one. The count is itself predictable from the AOI's WRS geometry: every region spanning two rows has more than one scene per date, and every region spanning one row has exactly one.

\subsection{Index normalisation: sea inside a scene-fitted dryness index}\label{sec:4.4}

TVDI is normalised against wet and dry edges taken as percentiles of the land surface temperatures a scene contains, within bins of a vegetation index. The AOIs are place-based rectangles that are not clipped to the coastline, and the water bit is deliberately preserved, so sea takes part in that fit. Water-dominant cells are 57.6 \% of one AOI and 38.9 \% of another, against 0.1 \% for the one inland region.

The contamination is real and visible: in the marine AOI the three lowest vegetation-index bins carry dry edges of 28.8 to 29.9 \textdegree{}C, which is sea-surface temperature and not a land dry edge, where the inland AOI's equivalent bins sit near 49 \textdegree{}C.

\textbf{It does not reach the modelled population.} Not one cell of the primary natural-vegetation population in any region has a mean vegetation index below 0.15, and the fifth percentile is 0.376 in the marine AOI, so the population occupies bins where the dry edge is a genuine land edge, 44 to 47.5 \textdegree{}C. Refitting the edges on land pixels alone shifts the signed association of the index with burning by an amount that scales with sea fraction, +0.017 and +0.016 in the two marine AOIs against +0.002 inland, and \textbf{changes no region's direction}.

A stronger test removes the scene dependence itself. Fitting one set of edges over the pooled land pixels of all five regions, 26.2 million of them, so that a given index value denotes the same dryness everywhere, leaves the disagreement intact: one region still ranks burned cells by higher dryness at 0.565 while two others rank them by lower dryness with bootstrap support, at 0.341 and 0.378. \textbf{The index's scene dependence is not the explanation for its reversal.}

\subsection{Label quality}\label{sec:4.5}

\textbf{A second burned-area product is not available for these years.} The usual comparator for MCD64A1 stops at December 2020 and every fire in this cohort is 2021 or 2022. That limitation cannot be closed with the products available and is stated rather than worked around.

\textbf{Agreement.} The share of a burned cell's positive sub-pixels agreeing with its modal burn date has a median of 1.000 in every region but a non-empty lower tail. Dropping burned cells below agreement thresholds of 0.75 and 0.90 leaves every reversal in the cohort intact, and moves the sharpest one \emph{further} from chance as thin-evidence cells are removed.

\textbf{Omission.} An independent active-fire observation covers the period and is independent of the burned-area algorithm. Cells it saw burning that the burned-area product did not label do sit systematically lower than labelled ones, 258 m against 512 m in one region, which is the predictor-correlated omission the literature worries about. Relabelling every such cell as burned, which deliberately over-corrects, changes \textbf{no} feature's side of chance in any region, and moves the sharpest reversal from 0.374 to 0.308, further from chance. The bias runs opposite to the direction that would threaten the result.

\textbf{Pre-label exposure.} Where the safeguard that removes cells burning inside their own predictor window did not run, the exposure is zero cells in one region and 28 in the other, under 0.2 \% of its population. It could not be measured from the archive, because both label rasters there are clipped to the label window; it required the unclipped product.

\subsection{Derived channels and design choices}\label{sec:4.6}

\textbf{Gap-filled thermal cells.} Where the fused product falls back on a modelled surface, the channel is not an observation. The gap-filled share is concentrated rather than diffuse, reaching 18.8 \% of cells in one region against 2.0 to 9.1 \% elsewhere. Restricting to cells at most 10 \% gap-filled leaves the increment bootstrap-supported in all five regions; the most exposed region moves most and in the predicted direction, +0.056 to +0.043, and the others by at most 0.008.

\textbf{Coordinate-bearing channels.} Two derived channels inherit a coordinate-derived component from their own model's inputs, which is the one route by which coordinates re-enter a feature set that excludes them. Dropping both retains 82 \% to 103 \% of the increment, every interval excluding zero, so the coordinate component is not what produces the within-region skill.

\textbf{Population definition.} The secondary all-valid population is \textbf{57.7 \% and 38.9 \% seawater} in the two marine AOIs, because sea cells satisfy every validity criterion. Most of the extra separability that population shows is therefore land against sea rather than any fire-relevant contrast, which is why the natural-vegetation population is primary.

\textbf{Population size and positive count.} Matching the largest region to the smallest on cell count and positive count together, over ten stratified draws, leaves its transfer behaviour inside the matched range in both roles while moving its within-region increment from +0.116 to a median +0.098. Part of a within-region increment is a positive-count effect; the transfer behaviour is not.

\textbf{Blocking scale.} Recomputing the transfer verdicts at 5 km rather than 1 km blocking moves the counts from ten positive, seven negative and three uncertain to six, four and ten, while \textbf{no direction changes side of the chance line and no point estimate changes sign}. Counts are the fragile quantity; signs are not.

\subsection{Software}\label{sec:4.7}

With byte-identical data, pipeline and seed, changing only the library version moves two frozen cross-region probes by +0.021 and +0.026 AUC, so cross-region point estimates carry an implementation tolerance of roughly $\pm$0.02 to 0.03 unless the version is pinned. Within-region AUCs reproduce to about four decimals across versions.

Against that, the within-region results rebuild \textbf{bit for bit}. Reconstructing the environment from the released dependency pins on a different operating system and re-running the released modelling step against the archived inputs reproduces every numeric field of the frozen metrics: 142 fields for one region with a maximum absolute difference of 6.9$\times$10$^{-18}$, and 168 for another with a maximum difference of \textbf{exactly 0}. A regularisation constant that the released sensitivity sweep omits was also recomputed directly and leaves every direction on its side of chance, widening the sweep's reported spread from 0.014 to 0.019.



\section{Discussion}\label{sec:5}

\subsection{The budget assembled}\label{sec:5.1}

Eight axes, measured on one cohort with everything else fixed, do not carry equal weight. Ordering them by how much they move a reported association gives a shape worth having in front of a reader (Table~\ref{tab:4}).

\begin{table}[htbp]
  \centering
  \small
  \caption{The budget assembled: eight axes ordered by how much they move a reported association. Every effect is measured on the same cohort with everything else held fixed, so the rows are comparable with each other and with the within-region effect these models report (+0.056 to +0.153 ROC-AUC, interval about $\pm$0.010 wide at 1 km blocking). The final column marks the axes whose relevance a reader can settle from their own inputs before fitting anything.}
  \label{tab:4}
  \begin{tabularx}{\linewidth}{>{\raggedright\arraybackslash}X>{\raggedright\arraybackslash}X>{\raggedright\arraybackslash}X}
    \hline
    Axis & Effect on the reported association & Predictable in advance? \\
    \hline
    Compositing chain & $\pm$0.02 AUC where it acts; \textbf{exactly zero where it cannot} & \textbf{yes}, from scenes per date \\
    Library version & $\pm$0.02 to 0.03 on cross-region estimates; $\sim$10$^{-4}$ within region & yes, from the pinned version \\
    Blocking scale & verdict counts move substantially; no sign changes & yes, by construction \\
    Index normalisation (sea) & +0.002 to +0.017 on a signed association; no direction changes & \textbf{yes}, from water fraction and index range \\
    Population size and positives & +0.018 on a within-region increment; none on transfer & yes, from population counts \\
    Gap-filled cells & up to $-$0.013 on an increment; support retained & yes, from the gap-fill share \\
    Quality screening (one region propagated) & below +0.0003 on a signed association, despite changing 92 \% of downscaled cells & no, it must be run \\
    Label omission & moves the sharpest association \emph{away} from chance & partly \\
    Coordinate-bearing channels & 0 to 18 \% of an increment & yes, from feature-importance records \\
    \hline
  \end{tabularx}
\end{table}

One axis is missing from that table because it is not a variation at all: the cell geometry, which is a description of what the grid is. The quality-screening split has now been propagated for one of the three unscreened regions, and its downstream effect on the reported association is \textbf{below +0.0003} despite changing 92 \% of that region's downscaled cells. It is left out of the table because one region does not fix a range, not because it is unmeasured.

\subsection{What should be reported}\label{sec:5.2}

The practical output of a budget is a reporting list. Six items cost nothing to report and would let a reader judge most of the above without repeating any of it.

\begin{enumerate}
  \item \textbf{Scenes per distinct acquisition date}, per region and window. This alone decides whether the compositing choice can matter.
  \item \textbf{Water-dominant fraction of the AOI}, and the range of the binning variable that the modelled population occupies, for any scene-fitted index.
  \item \textbf{The quality-screening rule actually applied to each input, per region}, with the date it was applied. A cohort split by export date is invisible otherwise.
  \item \textbf{The nodata convention}, explicitly. A product that declares none and writes zeros is indistinguishable in the output from one that observed a physical zero.
  \item \textbf{The analysis cell's ground dimensions in both axes}, not its nominal name, and the axis that block-size labels refer to.
  \item \textbf{The pinned library versions}, since the cross-region tolerance here exceeds several of the effects in the table above.
\end{enumerate}

\subsection{What generalises and what does not}\label{sec:5.3}

The magnitudes are this cohort's. Two things are not.

The \textbf{rules} of Section~\ref{sec:4.3} and \ref{sec:4.4} are statements about inputs. Any user can count scenes per date and any user can compare their modelled population's index range against the bins where water sits. Both answer "can this decision matter here" before any computation. In this cohort the compositing rule answered no for two regions of five, and the index rule answered no wherever the modelled population sits away from the water-contaminated bins.

The \textbf{failure modes} of Section~\ref{sec:4.2} are structural. A quality rule added to an export script partway through a project will split any cohort by export date, and nothing downstream records it. A nodata convention that writes zeros will encode sea as a physical temperature in any coastal AOI. Neither depends on this code being unusual; both depend only on the ordinary way such projects accumulate.

What does not generalise is the direction of any specific effect. We report, for instance, that correcting label omission sharpens the association rather than dissolving it. That is a fact about these five regions, and the opposite result elsewhere would not contradict it.

\subsection{Relation to the companion paper}\label{sec:5.4}

Every axis that bears on the companion paper's findings was tested against them, and in every case the finding survived: the sign reversal that carries that paper holds under land-only index edges, under a pooled common edge, under both label-agreement thresholds, under an over-corrected omission relabelling, and under population and positive-count matching. Two axes could not be pushed that far. The quality-screening split was a candidate for one region's unexplained behaviour, with a measured mechanism and the right sign, and propagating it through the whole chain removes it: no signed association moves by more than +0.0003. That closes a competing explanation for the companion paper without supplying an alternative. The cell geometry, meanwhile, means that paper's blocking argument is weaker in longitude than its labels imply.

That is the honest relationship between the two papers. This one does not rescue the other, and it was not run to. It reports what the observations are made of, and where that touches a claim, it says whether the claim holds.

\subsection{Limitations}\label{sec:5.5}

(i) \textbf{One pipeline.} Every magnitude here is measured on one implementation over one cohort. The rules and the failure modes are offered as general; the numbers are not.

(ii) \textbf{Five regions, and not a sample of anything.} They were selected for the companion paper's question, not to span a preprocessing design space, so the axes are exercised over whatever range those five happen to provide. The compositing axis, for instance, has three regions on one side of its dividing line and two on the other, which is enough to establish the rule and not enough to characterise its distribution.

(iii) \textbf{The screening axis is measured in one region of three.} Manavgat's chain was rebuilt end to end; Muğla and Bejís were not, so the axis has one region's downstream evidence rather than three. Muğla is the more informative of the two that remain, since its zero-fill share is 38.3 \% against Manavgat's 8.1 \%. The cell geometry remains a description rather than a variation: we cannot report what a truly square cell would have given.

(iv) \textbf{The independent fire observation is a different quantity.} Active-fire detections are not a burned-area product, the relabelling built from them is deliberately crude, and no claim is made that the relabelled cells burned. It probes the omission mechanism; it does not measure omission.

(v) \textbf{The compositing comparison is upstream of performance for four regions.} Where the seam verdict is uncertain or the intervention is inert, the released tool declines the downstream ROC-AUC comparison by design. The $\pm$0.02 tolerance therefore has one region's evidence behind it, while the rule that predicts whether the decision can act has all five.



\section{Conclusions}\label{sec:6}

A satellite fire-susceptibility result is decided in part before any model is fitted, and that part can be measured rather than gestured at. Holding data, model and validation fixed across five Mediterranean regions, eight preprocessing decisions were varied in turn. Two move a reported association by more than the interval width the headline result carries. Four do not threaten it, and one of those corrects in the direction opposite to the convenient one. Two could only be bounded.

The finding we would most like carried forward is not a magnitude. It is that two of these decisions announce their own relevance in the inputs. Counting scenes per distinct acquisition date says whether a compositing choice can move anything; comparing the modelled population's index range against the bins where water sits says whether sea contaminates a scene-fitted index. Both cost nothing and both answered no for half of our regions.

The failure modes we would most like guarded against are also not magnitudes. A quality rule added to an export script partway through a project splits a cohort by export date, and nothing downstream records it. A nodata convention that writes zeros encodes sea as a physical temperature. An analysis grid built by aggregating a fixed count of pixels in a geographic coordinate system is square in degrees and not on the ground. None of these requires an unusual pipeline; each requires only the ordinary way such projects accumulate.

Six things cost nothing to report and would let a reader judge most of a budget like this without repeating it: scenes per distinct date, the water fraction and the population's index range, the quality rule actually applied per region and when, the nodata convention, the analysis cell's ground dimensions in both axes, and the pinned library versions.


% ------------------------------------------------------------ declarations --
% Elsevier requires all four of these at submission. The two marked NEEDS
% AUTHOR INPUT cannot be written from the repository and must be completed
% before the manuscript is uploaded.

\section*{CRediT authorship contribution statement}

\textbf{Yunus Emre Cogurcu:} Conceptualization, Methodology, Formal analysis, Investigation, Writing -- original draft, Writing -- review and editing, Supervision. \textbf{Emrehan Metin:} Software, Data curation, Investigation, Validation, Writing -- review and editing.
% NEEDS AUTHOR INPUT: confirm this split with the co-author before submission.

\section*{Declaration of competing interest}

The authors declare that they have no known competing financial interests or
personal relationships that could have appeared to influence the work reported
in this paper.

\section*{Funding}

This work is an output of project 17506, supported by the Scientific Research
Projects (BAP) unit of \c{C}ukurova University.
%% Provisional wording, 2026-08-14. Confirm before submission: the project's
%% full title, the correct rendering of the unit's name in English, and whether
%% the funder requires a specific acknowledgement sentence.

\section*{Data availability}

All satellite inputs are public and are obtained through Google Earth Engine. The pipeline whose decisions are audited here, its configuration, the export metadata from which the provenance splits in Section 4.2 were read, and the frozen numeric outputs of every axis in Section 4 are publicly available; the repository, the commit of record and the licence are given in the data and code availability statement of the Methods. The scene inventories behind Section 4.3 and the pooled land-pixel edge fit behind Section 4.4 are reproducible from the released export scripts. No digital object identifier is minted and no archival deposit exists.

% ---------------------------------------------------------------- figures --
% Captions are maintained in ../figure_captions.tex and included verbatim.


\bibliographystyle{elsarticle-harv}
\bibliography{../REFERENCES}

\end{document}
